\section{The quadratic Chabauty calculation for Type 10}\label{app:type10}

Type 10 has a centered chain of length four on the first side, one
reflected pair of chains of length two and one reflected pair of
singletons on that side, and a reflected pair of chains of length three
and two reflected pairs of singletons on the second side. We prove that
its simultaneous rational center conditions force a collision. The
argument uses two genus-three covers, a normalized height at 13 and a
sieve on their common elliptic quotient. The sieve uses the full group
of that quotient. A full Mordell--Weil basis of either genus-three
Jacobian is unnecessary for the final step.

\subsection{Reduction to simultaneous center squares}

Let the positive centers of the first side be $u,w$, and let $v$ be the
positive center of the length-three chain pair on the second side. Put
$a=u^2$, $b=v^2$, $c=w^2$. For a chain of length $L$, write
$O_L=\{j-(L-1)/2:0\leq j<L\}$. The two remaining squared singleton
centers on the second side have power sums
\[
 p_k=\frac12\sum_{d\in O_4}d^{2k}
       +\sum_{d\in O_2}(u+d)^{2k}+c^k
       -\sum_{d\in O_3}(v+d)^{2k},\qquad 1\leq k\leq4.
\]
With $e_1=p_1$, $e_2=(p_1^2-p_2)/2$, the necessary recurrence relations
are $F=p_3-e_1p_2+e_2p_1=0$ and
$G=p_4-e_1p_3+e_2p_2=0$. Define
\begin{align*}
 D_0&=8a^2-48ab+4a+48b^2+24b-9,\\
 N_0&=48a^2b-8a^2-192ab^2-168ab+50a
                   +160b^3+352b^2+10b-27,\\
 J_0&=16a^2-32ab-8a+16b^2-8b+1,\\
 H_0&=24a^2-80ab-100a+80b^2+40b+81.
\end{align*}
Direct expansion gives
\[
 F=-\frac3{16}(N_0-2D_0c),\qquad
 \operatorname{Res}_a(N_0,2D_0)
       =-128(4b-25)(4b-9)^2(4b-1)^3.
\]
Thus $D_0=0$ forces $b=1/4,9/4,25/4$, each of which puts a
second-side chain vertex in the centered first-side chain. For a valid
configuration $D_0\ne0$ and $c=N_0/(2D_0)$. Substitution gives
\[
 G=\frac{(4b-9)(4b-1)J_0H_0}{64D_0}.
\]
The first two factors force the same collisions. Moreover
\[
 J_0=16\bigl(u^2-(v+1/2)^2\bigr)
              \bigl(u^2-(v-1/2)^2\bigr),
\]
so its vanishing makes chains of different colors intersect. Hence
$H_0=0$. A parametrization of this conic in $a,b$ is
\[
 d=3s^2-10s+10,\quad p=6s^2-5s+5,\quad
 q=3s^2+26s+10,\quad r=3s^2-10s+100,
\]
\[
 a=\frac{2p}{d},\qquad b=\frac q{4d},\qquad c=\frac r{4d}.
\]
The denominator $d$ has no rational zero. The exceptional parameters
$s=0,4,4/3$ are colliding configurations, as is the parameter at
infinity. With $z=3s-5$ define
\begin{equation}\label{eq:t10-q}
 q_0=z^2+5,\quad q_1=4z^2+30z+80,\quad
 q_2=z^2+36z+185,\quad q_3=z^2+275.
\end{equation}
The three square conditions become $a=q_1/q_0$, $4b=q_2/q_0$,
$4c=q_3/q_0$. In particular every valid finite configuration makes
all four $q_i(z)$ have the same nonzero rational squareclass.
The polynomial identities and all exceptional cases are checked by
\cert{check_sixth_mode.py} and \cert{check_sixth_center_lift.py}.

Write $z=A/B$ with coprime integers $A,B$, and homogenize $q_i$ to
$Q_i(A,B)$. All four quadratics are positive over the reals. Their
common squarefree part is therefore a positive squarefree integer
$\delta$. A prime dividing $\delta$ divides every $Q_i$, and
$Q_3-Q_0=270B^2$ implies that it lies in $\{2,3,5\}$. Such a prime
cannot divide $B$. If $A,B$ are odd, then $Q_0\equiv6\pmod8$ and
$Q_3\equiv4\pmod8$, so their squareclasses differ. The other
primitive parity cases give $Q_0\equiv1\pmod4$. This removes
$\delta=3,15$ and leaves
\begin{equation}\label{eq:t10-delta}
 \delta\in\{1,5\},\qquad q_i(z)=\delta c_i^2\quad(0\leq i\leq3).
\end{equation}
In fact $A$ is odd and $B$ is even, since the other mixed parity
would give $Q_3\equiv3$ or $7\pmod8$. This is a squareclass reduction
for the original center conditions, including $q_2$.

\subsection{The elliptic quotient and its full rational group}

Put $b_0^2=q_0q_3$. The genus-one quartic is birational to
\begin{equation}\label{eq:t10-E}
 E:\quad Y^2=X(X^2-560X+72900),\qquad
 X=2b_0+2z^2+280,\quad Y=2zX.
\end{equation}
For $X\ne0$, its inverse is $z=Y/(2X)$ and
$b_0=(X-2z^2-280)/2$. These maps extend over the smooth projective
models, so the points where the affine formulas have poles are retained.
Let $P=(180,900)$ and $T=(0,0)$. We use
\begin{equation}\label{eq:t10-full-group}
 E(\mathbf Q)=\mathbf Z P\oplus\langle T\rangle,\qquad 2T=0.
\end{equation}
Here is the arithmetic justification of completeness. The two-isogeny
descent images on $E$ and its partner
$Y^2=X(X^2+1120X+22000)$ are
\[
 \{1,5\},\qquad \{1,-5,-11,55\}.
\]
For each signed squarefree divisor $d$ of the relevant cubic coefficient,
the primitive covering quartic is tested on the two primitive local
patches modulo 256 and 3. The surviving classes are all attained by
explicit rational points. The isogeny index formula gives rank one.
Good reduction at 7 and 13 bounds the torsion order by four. There is
one nonzero rational two-torsion point, and a half of $T$ would have
$X=\pm270$, where the cubic takes a negative value. Thus the torsion
is precisely $\langle T\rangle$.

To establish that $P$ generates the free quotient, use $x/270=a/b$
and the duplication map
\[
 [a:b]\longmapsto
 [27(a^2-b^2)^2:4ab(27a^2-56ab+27b^2)].
\]
For coprime $a,b$, the coordinate gcd divides $11880$, and the maximum
of the coordinate absolute values lies between
$(121/48)\max(|a|,|b|)^4$ and $440\max(|a|,|b|)^4$.
Telescoping gives, for the doubled canonical height $h_c$ and the
naive height $h=\log\max(|a|,|b|)$,
\[
 h-\tfrac13\log(51840/11)\leq h_c
               \leq h+\tfrac13\log440.
\]
The two-isogeny class of $P$ is 5 and that of $T$ is 1, so the index
of $P$ in the free quotient is odd. If its absolute value were at
least three, a generator would have $h<\log24$. Exhaustion of
$|a|,b\leq24$ gives only the abscissae
$0,2/3,3/2,15/8,8/15$, corresponding to $T$, $\pm P$,
$\pm(P+T)$, $\pm2P$, and $\pm(2P+T)$. This contradicts a proper
odd index. The cutoff is certified by the integer inequality
$27\cdot440\cdot51840^9<24^{27}\cdot11^9$.
The rank and saturation arguments are recorded in
\cert{check_sixth_center_lift.py} and \cert{check_rank_one_generator.py}.

Write every elliptic image as $nP+\epsilon T$ with $\epsilon\in\{0,1\}$.
The common squareclass in \eqref{eq:t10-delta} is 1 for even $n$ and
5 for odd $n$. Good reductions at $13,37,59,113$, where the orders of
$P$ are $6,40,16,120$, respectively, give the necessary restrictions
\begin{equation}\label{eq:t10-mod240}
 \begin{array}{c|c}
 \epsilon& n\pmod{240}\\ \hline
 0&0,1,120\\
 1&0,120,239
 \end{array}.
\end{equation}
At undefined affine points the local test retains the point. Thus these
restrictions lose no candidate at infinity or at $X=0$. The remaining
small indices $0,\pm1$ give precisely the known collision configurations.

\subsection{Genus-three targets and rank hypotheses}

For $\delta=1,5$, let $X_\delta$ be the smooth projective curve
\[
 \delta a_0^2=q_1(z),\qquad b_0^2=q_0(z)q_3(z).
\]
Every finite simultaneous center lift maps to $X_\delta$. The quotients
by the independent signs are a conic, the elliptic quartic above and
\[
 C_\delta:\quad \delta y^2=q_0q_1q_3,\qquad y=a_0b_0.
\]
There are six simple branch points for the degree-four cover of the
$z$-line, so $X_\delta$ has genus three, and
$J(X_\delta)\sim_{\mathbf Q}E\times J(C_\delta)$.
The following conic parameters give the even-degree model
\begin{equation}\label{eq:t10-model}
 W^2=F_\delta(t)=(N_\delta^2+5D_\delta^2)
                         (N_\delta^2+275D_\delta^2),\qquad
 z=N_\delta/D_\delta,\quad a_0=M_\delta/D_\delta,\quad
 b_0=W/D_\delta^2,
\end{equation}
\[
\begin{array}{c|lll}
\delta&N_\delta&D_\delta&M_\delta\\ \hline
1&80-t^2&4t-30&2t^2-30t+160\\
5&25t^2-120t+80&10t^2-8&-60t^2+100t-48
\end{array}.
\]
They satisfy $\delta M_\delta^2=4N_\delta^2+30N_\delta D_\delta+
80D_\delta^2$. The degree-eight polynomials are squarefree and have
good ordinary reduction at 13.

A Richelot isogeny sends $J(C_\delta)$ to the Jacobian of
\[
 D_k:\quad v^2=kz(z-1)(z+5)(z^2-68z-275),\qquad k=-15/\delta.
\]
On the monic model $y^2=x(x-k)(x+5k)(x^2-68kx-275k^2)$,
full two-descent has étale algebra
$\mathbf Q^3\times\mathbf Q(\sqrt{159})$. The global norm kernel
unramified outside $\{2,3,5,11,19,53\}$ has dimension 25 over
$\mathbf F_2$. Its dimension after successive local conditions is
\[
\begin{array}{c|rrrrrrrr}
 k&\mathrm{norm}&\mathbf R&2&3&5&11&19&53\\ \hline
 -15&25&23&18&15&11&8&5&5\\
 -3 &25&23&18&15&11&8&6&5
\end{array}.
\]
The local image dimensions are $2,5,3,4,4,4,3$ at the seven listed
places. Explicit local divisor classes attain each dimension, so these
are complete local images. The required squareclass basis in
$\mathbf Q(\sqrt{159})$ follows from Minkowski's bound $<13$ and
norm relations among the primes over $2,3,5,11$. Its independence is
checked by local characters. This supplies unconditional Selmer dimension
five. Three dimensions are rational two-torsion, and the following
classes $A_k,B_k$ supply the other two dimensions. On the monic model
use
\[
\begin{array}{c|c|ll}
 k&A_k=[(x,y)-\infty]&u_{B_k}(x)&v_{B_k}(x)\\ \hline
 -15&(225/4,320625/32)&x^2+1020x+81225&-13590x-1564650\\
 -3 &(33,10692)&x^2+60x+3249&-1026x-40014
\end{array}.
\]
The divisor convention for $B_k$ is $[\operatorname{div}(u,v)-2\infty]$.
Exact divisibility $u\mid f_k-v^2$ and the independent Kummer classes
prove rank two. These calculations implement the full odd-degree descent
of Schaefer and Poonen--Schaefer~\cite{Schaefer1995,PoonenSchaefer}, and are reproduced by
\cert{check_richelot_two_selmer.py}. Thus $J(X_\delta)$ has rank three.

At 13, both genus-two reductions have 176 Jacobian points. For
$\eta=(dx/(2y),x\,dx/(2y))$, the rows of
$(\log A_k/13,\log B_k/13)$ modulo 13 are
\[
 \begin{array}{c|cc|c}
 k&\log A_k/13&\log B_k/13&\det\\ \hline
 -15&(3,1)&(1,11)&6\\
 -3 &(9,1)&(10,9)&6
 \end{array}.
\]
The exact calculation adds 176 times each divisor to an auxiliary
unramified degree-two divisor, stays in charts with unit leading
coefficients, and integrates in its two residue disks. The omitted tiny
integral terms have valuation at least $m-v_{13}(m)\geq2$ for $m\geq2$.
Since $176$ is prime to 13, the formal logarithm lattice is
$13\mathbf Z_{13}^2$. The determinant proves that the subgroup generated
by $A_k,B_k$ is 13-saturated. An independent reduction proof uses the
coefficient rows $(1,10),(0,1)$ for $k=-15$ and $(1,7),(1,3)$ for
$k=-3$ at 181 and 199. Each pair has nonzero determinant modulo 13.
The certificates are \cert{check_thirteen_logarithms.py} and
\cert{check_thirteen_saturation.py}.

For the elliptic generator, exact addition gives
$\log(P)/13\equiv1\pmod{13}$. The quotient-product and Richelot
isogenies have two-primary kernels. They therefore preserve logarithm
rank and saturation at 13. Both $J(X_\delta)(\mathbf Q)$ have logarithm
rank three. Their rational Néron--Severi rank is two. Indeed, the
Rosati-fixed fields of the genus-two reductions at 17 and 23 are
$\mathbf Q(\sqrt2)$ and $\mathbf Q(\sqrt{17})$, forcing the rational
Rosati-fixed algebra of $J(C_\delta)$ to be $\mathbf Q$. The two
nonisogenous factors of dimensions one and two give the additional
class. Consequently the quadratic Chabauty inequality is
$3<3+2-1$. We use the Mordell--Weil restricted height construction of
\cite{BalakrishnanDogra2018} and
make no assumption on the finiteness of the 13-primary Tate--Shafarevich
group.

\subsection{The correspondence and its normalized height}

Let $\iota(t,W)=(t,-W)$ and let the elliptic involution be
\[
 \sigma_1(t,W)=\left(\frac{15t-160}{2t-15},
                         \frac{9025W}{(2t-15)^4}\right),\qquad
 \sigma_5(t,W)=\left(\frac{24-25t}{25-30t},
                         \frac{9025W}{(25-30t)^4}\right).
\]
It fixes $z,b_0$ and changes the sign of $a_0$. The trace-zero
self-adjoint endomorphism is $Z=1+3\sigma$, with eigenvalues 4 and
$-2$ on the elliptic and genus-two factors. Put
\[
 B_1=(15/2,9025/16),\qquad B_5=(5/6,9025/144),\qquad
 H=\infty^++\infty^-.
\]
In both cases $\sigma(B_\delta)=\infty^+$. The four-point orbit of
$B_\delta$ under $\iota,\sigma$ is rational.

Use the proper graph correspondence
$\Gamma=\Gamma_\iota-3\Gamma_\sigma$. Its fiber degrees are $-2$,
and its diagonal intersection has degree $8-3\cdot4=-4$, giving
trace-zero degree after subtracting the two fibers. If $R_\tau$ denotes
the fixed divisor of an involution, the associated degree-zero divisor is
\begin{equation}\label{eq:t10-graphdiv}
 D_\Gamma(B,Q)=R_\iota-3R_\sigma-\iota Q+3\sigma Q
                                      -\iota B+3\sigma B.
\end{equation}
This divisor has class $Z[Q-B]+2Z[B-H/2]$. For
$j_H(Q)=[Q-H/2]$ and $\mathfrak q(A)=h_X(A,ZA)$, bilinearity gives
\[
 h_X([Q-B],D_\Gamma(B,Q))
      =\mathfrak q(j_H(Q))-\mathfrak q(j_H(B)).
\]
This is the centered height identity used below. The correspondence and finite local value theorem of
Balakrishnan--Dogra~\cite[Definition~6.2 and Theorem~6.3]{BalakrishnanDogra2018}
apply to this divisor.

Here all heights use $\log_{13}(13)=0$ and a compatible isotropic
complement to the holomorphic differentials. The elliptic complement is
the specified algebraic complement spanned by its second-kind quotient
form. The genus-two complement is transported through the Richelot
correspondence. Its pullback matrix on second-kind cohomology, in the
quotient basis, is
\[
 \begin{pmatrix}
 2&0&0&0\\0&2&0&0\\0&-32/3&-8/15&0\\
 -2200/9&-2464/9&-1114/45&-16/45
 \end{pmatrix}.
\]
Exact primitives verify the two nonholomorphic rows, and the cup matrices
satisfy the polarization identity with multiplier two. In particular the
transport fixes the height splitting, including its sign.
See \cert{check_richelot_de_rham_transport.py}.

For the two genus-two classes above, the global height matrices and the
scaled logarithm matrices are, modulo $13^3$ and $13^2$ respectively,
\[
\begin{array}{c|c|c}
 k&\mathcal H_k&L_k\\ \hline
 -15&\begin{pmatrix}923&1586\\1586&689\end{pmatrix}
     &\begin{pmatrix}68&92\\105&89\end{pmatrix}\\[3pt]
 -3&\begin{pmatrix}2093&676\\676&1027\end{pmatrix}
     &\begin{pmatrix}126&23\\105&22\end{pmatrix}
\end{array}.
\]
Thus $\mathcal A_k=L_k^{-t}(\mathcal H_k/13)L_k^{-1}$ is
\begin{equation}\label{eq:t10-heightmat}
 \mathcal A_{-15}=\begin{pmatrix}37&131\\131&135\end{pmatrix},\qquad
 \mathcal A_{-3}=\begin{pmatrix}90&114\\114&36\end{pmatrix}
                  \pmod{169}.
\end{equation}
The complete finite-prime sums for the three doubled mixed divisors
$(D_A,D_B)$, $(D_A,D_{A-B})$, $(D_B,D_{A-B})$ are
\[
\begin{array}{c|lll}
-15&-2\log(17)-2\log(593)&-2\log(5)&
                      2\log(7)+2\log(19)+2\log(313)\\
-3&-\log(3)&2\log(19)&-2\log(29)-2\log(31)
\end{array}.
\]
The dyadic terms vanish by smooth charts and signed intersection
cancellation. The finite terms depend on the chosen representatives.
For the shared-disk pairing on $D_{-3}$ one instead moves to
$4A-B$, computes its own finite intersections, and uses global
bilinearity. In that calculation the away sum is
$\log(3)-2\log(7)-2\log(61)+2\log(569)+2\log(1188727)$.
The raw Coleman integrals are corrected by their holomorphic de Rham
projections before adding these finite terms. The certificates
\cert{check_remaining_qr_heights13.py} and
\cert{check_complete_height_blocks13.py} reconstruct the displayed
matrices with explicit guard precision. This also records the factor
of two between a Jacobian class and its doubled anti-divisor.

In raw quotient logarithms $L=(L_e,L_{w,1},L_{w,2})$, write the
global quadratic expression as $L^t\mathcal J_kL/13$, with
\[
 \mathcal J_k=\operatorname{diag}(144,-\mathcal A_k/2).
\]
Adjunction and the Richelot polarization give
\[
 h_X(D,ZD)=2h_E(f_*D,f_*D)
                    -\tfrac12h_{D_k}(\phi g_*D,\phi g_*D).
\]
The elliptic coefficient satisfies
$13h_E(P,P)/\log_E(P)^2\equiv72\pmod{169}$.
The corresponding differential coordinates are
\[
 \frac{dz}{2b_0},\qquad \frac{dz}{ka_0b_0},\qquad
                              \frac{z\,dz}{a_0b_0}.
\]
Their pullback identities determine the normalization uniquely.

\subsection{The exact away-from-13 target sets}

Define the analytic function
\begin{equation}\label{eq:t10-rho}
 \rho_B(Q)=h_{13}(Q-B,D_\Gamma(B,Q))
               -\mathfrak q(j_H(Q))+\mathfrak q(j_H(B)).
\end{equation}
For rational $Q$ its value is the negative sum of the away local
heights. The complete target sets are
\begin{equation}\label{eq:t10-targets}
 \Omega_1=\{0,-1/6,-2/3\}\log_{13}(3),\qquad
 \Omega_5=\{0,1/2,2/3,-4/3\}\log_{13}(3).
\end{equation}
We give the local geometry that establishes exhaustiveness.

At every good prime other than 13, the normalized curve-level height
vanishes. The model discriminants restrict the remaining primes to
$2,3,5,11,19,53$. At 11 and 53 there is one geometric component with
ordinary nodes of thickness one. A local section cannot meet a node
with equation $xy=\ell$ over an unramified extension, since both
coordinates would have positive integral valuation. Every section
therefore retracts to the single vertex. At 19 the inner chart
$t=-2+19u$ for $\delta=1$, or $t=4+19u$ for $\delta=5$, and the
outer normalization give two elliptic components joined by two
thickness-one nodes. The involution $\sigma$ exchanges the components.
It commutes with $Z$, so the graph potential has the same value at
both vertices, and normalization makes that value zero.

At 5, pass to $\mathbf Q_5(\pi)$ with $\pi^2=5$ and set
$t=\pi u$, $W=25V$ on $X_1$. The reduction has nodes at $u=\pm1$,
each of thickness two. The main component is fixed and the two
exceptional components introduced by resolution are exchanged by
$\sigma$. A $\mathbf Q_5$ point of $X_1$ has
$v_\pi(u)=2v_5(t)-1$, so cannot reduce to either node and meets the
main component. For $X_5$, the isomorphism $a_1=\pi a_5$ gives
$u=a_5-2z/\pi$. Its conic equation forces $v_5(z)\geq1$ and
$a_5^2\equiv1\pmod5$, so all sections meet the two exceptional
components. Each twist has just one accessible orbit of components,
and its normalized local height is zero.

At 2, over any finite unramified extension, $t=2u$ on $X_1$ would give
$F_1(2u)/16\equiv3\pmod4$, which is impossible for a square. Thus
$s=1/t$ is integral for all local points, including infinity. Put
\[
 N_\infty=80s^2-1,\quad D_\infty=4s-30s^2,\quad
 S_\infty=N_\infty^2+140D_\infty^2,\quad
 V_\infty=(s^4W+S_\infty)/2.
\]
Then
\[
 V_\infty^2-S_\infty V_\infty=-18225D_\infty^4/4.
\]
Here $S_\infty\equiv1\pmod2$, so this equation splits into two
analytic disks, exchanged by $\iota$. In a semistable vertex set
containing their boundary points, further refinements within these
disks are genus-zero hanging trees. They have neither graph homology
nor positive-genus component contribution to the correspondence
measure. The potential is constant on each disk, and $\iota$
symmetry makes the constants equal. Normalization makes them zero.
The curve $X_5$ becomes $X_1$ over the unramified extension
$\mathbf Q_2(\sqrt5)$, preserving $\iota,\sigma,Z$. Vanishing thus
holds for both twists. This argument uses analytic disks and the
semistable vertex-set refinement theorem, without assuming a computed
full wild regular model.

At 3, the cluster picture has a four-root cluster near $-1$ of depth
$1/2$ containing two twins of depth $5/2$, and a four-root cluster near
4 of depth $3/2$ containing a twin of depth 2. After the quadratic
ramified base change, the charts reduce to
$V^2=(u^2-1)^2$ and $V^2=u^2(u^2-1)$. The cluster-to-graph rules give
two edges $A^+\to A^-$ of length 4, bridges $A^+\to B$ and
$B\to A^-$ of length 2, and one loop at $B$ of length 1, with lengths
normalized by $v_3(3)=1$. The involution exchanges the long edges,
fixes the bridges and reverses the loop. Cycle vectors
$(1,-1,0,0,0)$, $(1,1,-2,-2,0)$, $(0,0,0,0,1)$ have weighted
Gram matrix $\operatorname{diag}(8,24,1)$, while $Z$ acts by
$(-2,4,-2)$. The edge densities of the graph potential are
$(-1/12,-1/12,2/3,2/3,-2)$.

Normalized at $B_1$, its values along a bridge, a long edge and the
loop are respectively
\[
 -\frac1{12}+\frac s3-\frac{s^2}3,\qquad
 \frac{x^2}{24}-\frac x6-\frac1{12},\qquad
 x^2-x-\frac34.
\]
Integer valuations leave accessible only $A^\pm$, bridge coordinates
$1/2,3/2$, and the loop midpoint. The midpoint is absent for $X_1$,
since $F_1(4+9u)/3^7\equiv2u^2+2u+1\pmod3$ has no residue root
and has odd valuation. Local witnesses attain the other positions.
For $X_5$ all four positions occur. After its basepoint normalization,
the exact potential sets are
\[
 \{-1/12,0,-1/3\},\qquad \{1/4,1/3,0,-2/3\}.
\]
An intersection calculation calibrates the scalar and sign. At $t=-1$
on $X_1$, choose $W/3\equiv-1\pmod3$ and take the opposite divisor
$-D_\Gamma$. Its principal term has valuation 2, its horizontal
intersection is $3/2$, and its vertical graph correction is $1/3$.
Its height is $(2-3/2-1/3)\chi_3(3)=-\log_{13}(3)/6$.
Returning to $D_\Gamma$ gives scalar $-2\log_{13}(3)$ for the
displayed potential. This gives \eqref{eq:t10-targets} with the
negative-away convention in \eqref{eq:t10-rho}.
The graph potential is given by \cite[Theorem~3.2]{BDMTV}. The
cluster construction uses \cite[Theorems~7.1 and~8.6]{DDMM}, and the
vertex refinement uses \cite[Proposition~3.13(3)]{BPR}. All polynomial charts, accessibility witnesses, graph
Laplacians and the scalar calibration are checked by
\cert{check_qc_normalization_and_local_targets.py} and
\cert{check_complete_away_targets.py}.

\subsection{Moving normalization and its precision}

Use the monic model $Y^2=f(t)=F_\delta(t)/c^2$, with $c=1,5625$
for $\delta=1,5$, and put $\omega_j=t^jdt/(2Y)$, $\rho=\omega_3$.
If $r_j$ is the coefficient of $u^j$ in
$(u^8f(1/u))^{-1/2}$, a residue-zero basis is
\[
 \eta=(\omega_0,\omega_1,\omega_2,
             \omega_4-r_1\rho,\omega_5-r_2\rho,\omega_6-r_3\rho).
\]
The triples $(r_1,r_2,r_3)$ are $(0,-2080,33600)$ and
$(208/75,6752/1875,-386816/84375)$. The compatible complement is Lagrangian. Its lower-block determinants
are $3/760$ and $4271484375/9728$. Both are 13-adic units, so holomorphic projection
loses no precision.

For the Coleman--Gross de Rham projection $\Psi$ and holomorphic
projection $\lambda$ along this complement, the infinity input is
\[
\begin{array}{c|c|c}
\delta&\Psi\rho\pmod{2197}&\lambda\rho\pmod{2197}\\ \hline
1&(1433,392,1539,1586,728,1248)&(289,925,590)\\
5&(1194,1844,368,1690,533,1898)&(427,1168,628)
\end{array}.
\]
To compute it, expand Frobenius with
$R=f(t^{13})-f(t)^{13}$ as
\[
 13t^{13(i+1)-1}\sum_{n\geq0}\binom{-1/2}{n}
           R^n\frac{dt}{2Y^{13+26n}}.
\]
The retained terms $0\leq n\leq4$ are reduced modulo $13^{25}$ with
an auxiliary scale $13^8$. At most seven denominator divisions lose
digits. For the omitted terms, removing finite branch principal parts
loses at most $\lfloor\log_{13}(11+26n)\rfloor$ digits and horizontal
reduction loses at most one. The resulting valuation bound
\[
 n+1-\lfloor\log_{13}(11+26n)\rfloor-1\geq4\quad(n\geq5)
\]
tends to infinity. The omitted infinite tail is therefore
$O(13^4)$. The residue row of affine Frobenius is
$(0,0,0,0,0,0,13)$. Its proper block $F_{\rm dR}$ and last column $v$
satisfy $(F_{\rm dR}-13I)\Psi\rho=v$. The determinant has valuation
three, but the inverse entries lose at most one digit, so precision
$13^4$ gives $\Psi\rho$ modulo $13^3$. This is the even-degree
Frobenius projection formula of \cite[Proposition~5.3]{GajovicMuller}. The complete exact implementation is
\cert{check_infinity_projection.py}.

For a finite source $Q=(a,y_a)$ use the third-kind kernel
\[
 \nu_Q=\frac12\left(1+\frac{y_a}{Y}\right)\frac{dt}{t-a},
\]
whose residue divisor is $Q-H/2$. At a source tending to infinity set
$h=1/a$, $v=h^4y_a$ and use instead
\[
 \nu_\infty(h,v)=-\frac{h\,dt}{2(1-ht)}
                       -\frac{v t^3dt}{2Y(1-ht)}.
\]
It differs from the finite kernel by a holomorphic differential and
has the same normalized third-kind form. It specializes to
$-\rho$ at $\infty^+$.
For a moving finite source, differentiation modulo an exact form gives
\[
 \partial_a\nu_Q\equiv
 \frac{(f'(a)+f'(t))(t-a)-2(f(t)-f(a))}{(t-a)^2}
                       \frac{dt}{4y_aY}.
\]
Its numerator quotient is polynomial of degree at most six. The
analogous expression in the infinity chart is also polynomial and
residue zero. The projection commutes with this differentiation by
its global-symbol characterization: differentiate the analytic
Coleman primitives against fixed second-kind test forms and take
residues coefficientwise. Nondegeneracy of the cup product then
identifies the derivative of the projection.

A residue-correct differential for \eqref{eq:t10-graphdiv} is
\[
 d\log(W/M^3)+3\nu_{\sigma B}-\nu_{\iota B}
                         +3\nu_{\sigma Q}-\nu_{\iota Q}.
\]
Subtracting its holomorphic projection defines the local height.
For $q=t(Q)-t(B)$, integration on the basepoint disks yields
\[
 h_{13,1}=991q+2083q^2,\qquad
 h_{13,5}=1627q+81q^2\pmod{2197}.
\]
For the opposite-sheet kernel one uses its nonsingular expression
\[
 \nu_{\iota Q}=
 \frac{f(t)-f(t(Q))}{t-t(Q)}
              \frac{dt}{2Y(t)(Y(t)+Y(Q))}.
\]
All denominators are units on the relevant bidisk. The integral of a
term of total order $n\geq3$ has valuation at least
$n-v_{13}(n)\geq3$. For the twice-integrated moving normalizer the
bound is $n-v_{13}(j)-v_{13}(i+1)\geq3$, where $n=i+j+1$.
These estimates tend to infinity and control analytic tails.

The centered basepoint logarithms in the three quotient coordinates
are $(0,1313,832)$ for $\delta=1$ and $(1339,1339,1079)$ for
$\delta=5$, modulo 2197. They are determined from the signed cup
matrix and $\Psi\nu_B$, with
$\langle\Psi(\partial_a\nu_Q),\omega_j\rangle=-a^j/(2y_a)$.
If the raw logarithm increment is $d_1q+d_2q^2+\cdots$, the global
change through precision $13^3$ is
\[
 \frac{2L_B^t\mathcal J_kd_1q+
 (d_1^t\mathcal J_kd_1+2L_B^t\mathcal J_kd_2)q^2+
 2d_1^t\mathcal J_kd_2q^3}{13}.
\]
The cubic cross term is retained. Every omitted product of increment
degrees $a,b$ with $a+b\geq4$ has valuation at least
$a+b-v_{13}(a)-v_{13}(b)-1\geq3$.
Putting $q=13u$ now gives
\begin{align}\label{eq:t10-rhoseries}
 \rho_{B_1}/13&=149u+136u^2\pmod{169},\\
 \rho_{B_5}/13&=129u+110u^2+117u^3\pmod{169}.\nonumber
\end{align}
The entire analytic functions have integral coefficients on
$\mathbf Z_{13}$. All terms of degree at least three vanish modulo
13. Thus their simple residue roots exhaust the exact analytic roots,
by Hensel's lemma. The coefficients and both tail estimates are part
of the proof, while \cert{check_sixth_mode_qc_roots13.py} checks the
finite algebra and the congruences.

\subsection{The 48 local roots and the final elliptic sieve}

The mod-240 condition and squareclass parity retain four of the twelve
residue disks on each curve. On $X_1$ their reductions are
$(1,\pm1)$ and the two infinity branches with $W/t^4=\pm1$.
On $X_5$ they are $(3,\pm3)$ and the infinity branches with
$W/t^4=\pm9$. Each set is the free orbit of its basepoint disk under
$\{1,\iota,\sigma,\iota\sigma\}$.
Since $\log_{13}(3)/13\equiv5\pmod{13}$, reducing the exact target
equations from \eqref{eq:t10-targets} gives
\[
\begin{array}{c|c|c|c}
\delta&\text{target}/\log_{13}(3)&u\pmod{13}&
        \text{derivatives modulo }13\\ \hline
1&0&0,12&6,7\\
1&-1/6&2,10&4,9\\
1&-2/3&\varnothing&\\ \hline
5&0&0,11&12,1\\
5&1/2&2,9&10,3\\
5&2/3&3,8&9,4\\
5&-4/3&5,6&7,6
\end{array}.
\]
Every displayed residue gives exactly one exact local root. Thus there
are four roots on the $B_1$ disk and eight on the $B_5$ disk.

For completeness of transport, put
$Z_{\rm act}=-\iota_*+3\sigma_*$. Direct subtraction of correspondence
divisors gives
\[
 D_\Gamma(A,Q)-D_\Gamma(B,Q)=Z_{\rm act}(A-B),\qquad
 D_\Gamma(B,Q)-D_\Gamma(B,A)=Z_{\rm act}(Q-A).
\]
Local-height symmetry and self-adjointness cancel the cross terms,
so $h_A(Q)=h_B(Q)-h_B(A)$. Equivariance under $\iota,\sigma$
then shows that $h_B(\tau Q)-h_B(Q)$ is constant. Applying
$\tau$ twice makes it zero. The argument is first made with all
supports disjoint and extends by the normalized analytic kernel.
The centered global term has the same invariance. Transporting the
base-disk equations therefore gives exactly $16+32=48$ local roots
on the eight retained disks. Eight are the known basepoint-orbit
roots. The other forty are local roots to be tested, without an
assertion that they are rational points.

At each of the good primes $61,101,107$, enumerate the complete cycle
of $P$ and its translate by $T$. Their orders are $26,13,104$.
At a finite point with $X\ne0$, form $z=Y/(2X)$ and require the four
$q_i(z)$ in \eqref{eq:t10-q} to have the same nonzero quadratic
character. A zero value, the point at infinity, and $X=0$ are retained
conservatively. Projecting each permitted set to the gcd of the point
order with 3120 and combining with \eqref{eq:t10-mod240} leaves
\[
\begin{array}{c|c}
\epsilon&n\pmod{3120}\\ \hline
0&0,1,1560\\
1&0,1560,3119
\end{array}.
\]
This finite group calculation is used on the isolated height roots.
On the $B_1$ disk the elliptic reduction is the identity, the torsion
index is $\epsilon=0$, and $n$ is even. Hence the sieve forces
$n\equiv0\pmod{13}$. On the $B_5$ disk the reduction is $-P+T$,
$\epsilon=1$ and $n$ is odd, giving $n\equiv12\pmod{13}$.

The elliptic quotient differential is $dz/(2b_0)$ and its
$t$-derivative at both basepoints is 11 modulo 13. With
$\log(P)/13\equiv1\pmod{13}$, this gives
\[
 n\equiv11u\pmod{13}\quad(B_1),\qquad
 n\equiv12+11u\pmod{13}\quad(B_5).
\]
The logarithm calibration is independently checked by exact rational
addition: the formal parameter of $6P$ is $78\pmod{169}$, and its
formal logarithm is the same modulo 169. Only $u=0$ survives on
either basepoint disk. The exact zero-target equation has the root
$B_\delta$ there, and simple-root uniqueness excludes a second root
in that subdisk. The original center-square conditions are invariant
under both involutions, so the same conclusion holds on every
retained disk. All simultaneous center-lift candidates are among the
eight rational orbit points. Their elliptic indices are $0,\pm1$,
and the exact collision tests exclude them. This proves the Type 10
exclusion at every rational height.

\subsection{Arithmetic dependencies and reproducibility}

The numerical inputs in this appendix are obtained from different
certificates. Their roles are summarized here to make the proof
reproducible from the stated models.
\begin{longtable}{@{}p{.23\textwidth}p{.70\textwidth}@{}}
\toprule
Input & Certificate and mathematical role\\ \midrule
Center reduction & \cert{check_sixth_mode.py},
\cert{check_sixth_center_lift.py},
\cert{check_sixth_genus_two_covers.py}. Moment elimination,
exceptional collision parameters and coverage by $\delta=1,5$.\\
Elliptic group & \cert{check_rank_one_generator.py},
\cert{check_lift_congruences.py}. The height cutoff proves full
generation, and finite reductions give the mod-240 conditions.\\
Genus-two ranks & \cert{check_richelot_two_selmer.py}. Full local
images and an explicit global squareclass basis prove rank two
unconditionally.\\
Logarithms and index & \cert{check_thirteen_logarithms.py},
\cert{check_thirteen_saturation.py}. Unit normalized determinant and
independent finite-field tests prove full logarithm rank and
13-saturation.\\
Compatible splitting & \cert{check_genus_three_qc_de_rham.py},
\cert{check_richelot_de_rham_transport.py}. Exact cup products,
correspondence action and full second-kind pullbacks fix conventions.\\
Global heights & \cert{check_elliptic_height13.py},
\cert{check_complete_height_blocks13.py},
\cert{check_remaining_qr_heights13.py}. Normalized Coleman integrals
plus the complete representative-specific finite terms give
\eqref{eq:t10-heightmat}.\\
Curve-level local values & \cert{check_qc_normalization_and_local_targets.py},
\cert{check_complete_away_targets.py}. Regular charts, metric graph
and intersection calibration give exactly \eqref{eq:t10-targets}.\\
Infinity projection & \cert{check_infinity_projection.py}. Affine
Frobenius and the proved tail bound give $\Psi\rho$ and $\lambda\rho$.\\
Moving roots & \cert{check_sixth_mode_qc_roots13.py}. It recomputes
the derivative reductions, normalized local heights and centered
quadratic term, and checks simple-root residue lists.\\
Final candidate sieve & \cert{check_sixth_mode_candidate_sieve.py}.
Complete elliptic cycles at three primes and independent logarithm
calibration remove every non-basepoint local height root.\\ \bottomrule
\end{longtable}

The root certificate uses the global height matrices, the infinity
projection and the rank statements as inputs established by the earlier
certificates. It verifies the root series using the coefficients it
actually computes. Its Hensel refinements modulo 169 are reproducibility
outputs. The infinite-tail bounds, metric-graph interpretation and
height-cocycle argument are mathematical parts of the proof.
The final sieve requires no saturation at untested odd primes of the
genus-three Jacobians because it uses \eqref{eq:t10-full-group}.
The conclusion classifies the simultaneous-center-lift candidates
needed by Type 10. The full rational-point sets $X_1(\mathbf Q)$ and
$X_5(\mathbf Q)$ can contain points failing the remaining center-square
conditions and are outside this assertion.
